%% kapa: exam=1 round=32 year=2021 period=2 form=A subject=accounting subject_ko=회계학 seq=33 number=73
%% source: Q-Net 감정평가사 기출문제 (cst00309) artlSeq=5208397 file=extracted/1008/5208397/A형/2021년도 제32회 감정평가사 1차 2교시 A형.pdf
%% conversion: kordoc 4.24.51 -> tools/gen1.py (auto); review=proofread-sonnet; flags=none
\begin{question}{73}
(주)감평은 제품라인인 A, B, C부문을 유지하고 있다. 20×1년 각 부문별 손익계산서는 다음과 같다.

{\small\centering\begin{tabular}{|l|r|r|r|r|}\hline
 & \multicolumn{1}{c|}{A부문} & \multicolumn{1}{c|}{B부문} & \multicolumn{1}{c|}{C부문} & \multicolumn{1}{c|}{합계} \\ \hline
매출액 & ￦200,000 & ￦300,000 & ￦500,000 & ￦1,000,000 \\
변동원가 & 100,000 & 200,000 & 220,000 & 520,000 \\ \hline
공헌이익 & 100,000 & 100,000 & 280,000 & 480,000 \\
고정원가 & & & & \\
\hspace{1em}급여 & 30,000 & 50,000 & 80,000 & 160,000 \\
\hspace{1em}광고선전비 & 10,000 & 60,000 & 70,000 & 140,000 \\
\hspace{1em}기타 배부액 & 20,000 & 30,000 & 50,000 & 100,000 \\ \hline
영업손익 & ￦40,000 & (￦40,000) & ￦80,000 & ￦80,000 \\ \hline
\end{tabular}\par}
\vspace{4pt}
(주)감평의 경영자는 B부문의 폐쇄를 결정하기 위하여 각 부문에 관한 자료를 수집한 결과 다음과 같이 나타났다.
\begin{framedtext}
○ 급여는 회피불능원가이다. \par
○ 광고선전은 각 부문별로 이루어지기 때문에 B부문을 폐쇄할 경우 B부문의 광고선전비는 더 이상 발생하지 않는다. \par
○ 기타 배부액 총 ￦100,000은 각 부문의 매출액에 비례하여 배부한 원가이다. \par
○ B부문을 폐쇄할 경우 C부문의 매출액이 20\%  감소한다. \par
\end{framedtext}
(주)감평이 B부문을 폐쇄할 경우 (주)감평 전체 이익의 감소액은? (단, 재고자산은 없다.)
\begin{choices}
  \item ￦36,000
  \item ￦46,000
  \item ￦66,000
  \item ￦86,000
  \item ￦96,000
\end{choices}
\end{question}
\answer{5}
